% ==========================================
% BAB V RENCANA SELANJUTNYA
% ==========================================
\chapter{RENCANA SELANJUTNYA}
\label{chap:rencana-selanjutnya}
\section{Rencana Implementasi}

Rencana implementasi yang akan dibahas pada Bab ini berisikan alat dan bahan, 
\textit{timeline}, dan lingkungan dan konfigurasi. Berikut merupakan detail dari 
setiap subbab yang akan dibahas.

\subsection{Alat dan Bahan}

Berikut merupakan pembahasan dari perencanaan alat dan bahan yang digunakan
pada Tabel V.1 berikut.

\begin{longtable}{|p{1cm}|p{3.5cm}|p{8cm}|}
\caption{Daftar Alat dan Bahan Implementasi Sistem}
\label{tab:alatbahan} \\
\hline
\textbf{No} & \textbf{Layanan} & \textbf{Alat dan Bahan} \\
\hline
\endfirsthead

\hline
\textbf{No} & \textbf{Layanan} & \textbf{Alat dan Bahan} \\
\hline
\endhead

% Row 1 -----------------------------------------------------------------
1 &
\textit{Frontend} &
1. Framework: Next.js \newline
2. Bahasa: TypeScript \newline
3. Styling: Tailwind CSS \newline
4. UI Library (opsional): ShadCN / Chakra UI
\\
\hline

% Row 2 -----------------------------------------------------------------
2 &
\textit{Backend \& LLM Service} &
1. Node.js / Express untuk API \newline
2. OpenAI API (GPT-4o) untuk generasi konten \newline
3. LangChain untuk \textit{orchestration} dan \textit{prompt flow}
\\
\hline

% Row 3 -----------------------------------------------------------------
3 &
\textit{Database \& Storage} &
1. PostgreSQL untuk data pengguna dan progres belajar \newline
2. Supabase / Firebase Storage untuk penyimpanan dokumen materi
\\
\hline

% Row 4 -----------------------------------------------------------------
4 &
\textit{Infrastructure \& Deployment} &
1. Platform Hosting: Vercel / Netlify (Frontend) \newline
2. Railway (Backend) \newline
3. GitHub untuk \textit{version control}
\\
\hline

% Row 5 -----------------------------------------------------------------
5 &
\textit{Documentation \& IDE} &
1. Visual Studio Code \newline
2. \LaTeX{} (VS Code TexTools) \newline
3. Postman untuk uji API
\\
\hline

% Row 6 -----------------------------------------------------------------
6 &
\textit{Testing Tools} &
1. Black Box Testing Sheet untuk pengujian fungsionalitas sistem \newline
2. Google Form untuk \textit{usability testing} dan pengumpulan umpan balik
\\
\hline

\end{longtable}

\subsection{\textit{Timeline}}

\textit{Timeline} implementasi disusun dengan membagi proses pengembangan 
ke dalam beberapa fase yang terstruktur. Berikut adalah \textit{timeline} 
implementasi aplikasi pembelajaran ASD berbasis LLM.

\textbf{1. Fase 1: Perancangan (Oktober 2025 – Desember 2025)}
\begin{enumerate}[label=\alph*.]
    \item Identifikasi dan analisis permasalahan pembelajaran ASD
    \item Analisis kebutuhan pengguna serta penyusunan kebutuhan fungsional dan nonfungsional
    \item Studi literatur terkait LLM, \textit{guided learning}, dan pembelajaran ASD
    \item Analisis alternatif solusi dan pemilihan teknologi
    \item Perancangan desain konsep solusi, alur \textit{guided learning}, dan arsitektur sistem
\end{enumerate}

\textbf{2. Fase 2: Initial Setup dan Implementasi Fitur Dasar (Januari 2026 – Februari 2026)}
\begin{enumerate}[label=\alph*.]
    \item \textit{Setup} lingkungan pengembangan, repository, dan struktur proyek
    \item Perancangan arsitektur sistem secara detail
    \item Perancangan \textit{database} sistem
    \item Implementasi \textit{database} dan integrasi awal \textit{backend}
    \item Implementasi modul \textit{backend} dasar (autentikasi, materi, progres belajar)
    \item Implementasi UI/UX awal untuk halaman beranda dan pemilihan topik
    \item Integrasi awal API LLM untuk menghasilkan konten dasar
\end{enumerate}

\textbf{3. Fase 3: Implementasi \textit{Guided Learning} (Maret 2026 – April 2026)}
\begin{enumerate}[label=\alph*.]
    \item Penyusunan \textit{structured prompt flow} (konsep → contoh → latihan → ringkasan)
    \item Implementasi halaman materi berbasis \textit{guided learning}
    \item Integrasi LLM untuk menghasilkan konten sesuai kurikulum
    \item \textit{Service integration testing} untuk memastikan stabilitas alur sistem
\end{enumerate}

\textbf{4. Fase 4: Pengujian dan Evaluasi (Mei 2026 – Juni 2026)}
\begin{enumerate}[label=\alph*.]
    \item \textit{Black Box Testing} untuk seluruh fitur aplikasi
    \item \textit{Task-based usability testing} menggunakan Google Form
    \item Evaluasi pengalaman pengguna dan konsistensi materi
    \item Penyempurnaan fitur berdasarkan hasil evaluasi
    \item Finalisasi implementasi sistem dan dokumentasi hasil pengujian
\end{enumerate}

\subsection{Lingkungan dan Konfigurasi}

\begin{enumerate}
    \item Lingkungan pengembangan (\textit{development environment}) dilakukan 
    pada perangkat lokal menggunakan macOS dengan konfigurasi pendukung seperti 
    Node.js, Python, dan Docker untuk menjalankan layanan \textit{backend}, 
    \textit{frontend}. Seluruh proses dilakukan secara lokal 
    sebelum dipindahkan ke tahap integrasi.

    \item Lingkungan uji coba (\textit{testing environment}) dijalankan pada 
    \textit{staging server} menggunakan layanan seperti Vercel (frontend) dan 
    Railway (backend). Lingkungan ini digunakan untuk menguji alur 
    \textit{guided learning} dan integrasi LLM sebelum 
    aplikasi dinilai siap untuk demonstrasi dan evaluasi pengguna.
\end{enumerate}

\section{Desain Pengujian dan Evaluasi}

Pengujian sistem dirancang untuk memastikan bahwa seluruh fungsionalitas aplikasi
dan kualitas alur pembelajaran berjalan sesuai tujuan yang telah ditetapkan.
Berikut adalah desain pengujian dan evaluasi yang akan dilakukan:

\begin{longtable}{|p{3.2cm}|p{2cm}|p{3.5cm}|p{3.5cm}|}
\caption{Desain Pengujian dan Evaluasi Sistem}
\label{tab:pengujian-evaluasi}\\
\hline
\multicolumn{4}{|c|}{\textbf{Desain Pengujian dan Evaluasi Sistem}} \\
\hline
\textbf{Jenis Pengujian} & \textbf{Metode} & \textbf{Tujuan} & \textbf{Indikator Keberhasilan} \\
\hline
\endfirsthead

\multicolumn{4}{c}{\small (lanjutan dari halaman sebelumnya)}\\[0.3em]
\hline
\textbf{Jenis Pengujian} & \textbf{Metode} & \textbf{Tujuan} & \textbf{Indikator Keberhasilan} \\
\hline
\endhead

\hline
\multicolumn{4}{r}{\textit{Bersambung ke halaman berikutnya}}\\
\endfoot

\hline
\endlastfoot
\caption{Desain Pengujian dan Evaluasi Sistem}
\label{tab:pengujian-evaluasi}\\

\hline
\textbf{Jenis Pengujian} & \textbf{Metode} & \textbf{Tujuan} & \textbf{Indikator Keberhasilan} \\
\hline
\endfirsthead

\multicolumn{4}{c}{\small (lanjutan dari halaman sebelumnya)}\\[0.3em]
\hline
\textbf{Jenis Pengujian} & \textbf{Metode} & \textbf{Tujuan} & \textbf{Indikator Keberhasilan} \\
\hline
\endhead

\hline
\multicolumn{4}{r}{\textit{Bersambung ke halaman berikutnya}}\\
\endfoot

\hline
\endlastfoot

\textit{Black Box Testing} &
Verifikasi fungsional &
Memastikan seluruh fitur berjalan sesuai spesifikasi &
Semua skenario uji pada \textit{black box testing} lulus tanpa error fungsional \\
\hline

\textit{Availability Test} &
Simulasi akses berulang &
Memastikan aplikasi dapat digunakan tanpa gangguan selama sesi belajar &
Aplikasi stabil, tidak terjadi \textit{crash} atau \textit{downtime} selama pengujian \\
\hline

\textit{Security Validation} &
Pemeriksaan autentikasi dan validasi \textit{input} &
Memastikan hanya pengguna sah yang dapat mengakses sistem serta \textit{input} tidak valid ditolak &
Proses \textit{login} tidak dapat dilewati, percobaan kredensial salah dan \textit{input} tidak valid ditolak oleh sistem \\
\hline

\textit{Compatibility Test} &
Uji lintas perangkat dan \textit{browser} &
Memastikan aplikasi berjalan dengan baik di macOS, Windows, perangkat \textit{mobile}, serta \textit{browser} modern &
Halaman dan fitur utama tampil responsif dan berfungsi di Chrome, Firefox, Safari, serta \textit{browser} modern lain yang diuji \\
\hline

Pengujian Alur Aplikasi &
\textit{Task-based usability test} &
Memastikan alur materi dan tahapan \textit{guided learning} mudah diikuti pengguna &
Mayoritas peserta dapat menyelesaikan tugas tanpa kebingungan signifikan, dengan waktu penyelesaian tugas berada dalam rentang yang stabil \\
\hline

\end{longtable}



\section{Analisis Risiko dan Mitigasi}

Analisis risiko dan mitigasi disusun untuk mengidentifikasi potensi hambatan selama pengembangan serta menentukan strategi penanganan yang tepat agar proyek tetap berjalan sesuai rencana. Berikut adalah analisis risiko dan mitigasi:

\begin{longtable}{|p{3.5cm}|p{2cm}|p{1.5cm}|p{5.5cm}|}
\caption{Analisis Risiko dan Mitigasi Sistem Pembelajaran ASD Berbasis LLM} \\
\hline
\textbf{Risiko} & \textbf{Probabilitas} & \textbf{Dampak} & \textbf{Mitigasi} \\
\hline
\endfirsthead

\multicolumn{4}{c}{\textit{(Lanjutan dari halaman sebelumnya)}} \\
\hline
\textbf{Risiko} & \textbf{Probabilitas} & \textbf{Dampak} & \textbf{Mitigasi} \\
\hline
\endhead

\hline \multicolumn{4}{r}{\textit{Bersambung ke halaman berikutnya}} \\
\endfoot

\hline
\endlastfoot

1. LLM menghasilkan informasi tidak akurat (hallucination) 
& Sedang 
& Tinggi 
& Membatasi ruang lingkup \textit{prompt} menggunakan \textit{prompt template} yang spesifik, dan melakukan \textit{expert review} untuk memastikan akurasi konsep. \\
\hline

2. Keluaran LLM tidak sesuai kurikulum / tahap \textit{guided learning}
& Sedang
& Tinggi
& Menggunakan \textit{prompt template} spesifik, menerapkan \textit{guardrails}, dan memastikan seluruh konteks berasal dari \textit{Knowledge Base} ASD. \\
\hline

3. Gangguan kerja karena beban magang sehingga \textit{deadline} TA mundur
& Tinggi
& Tinggi
& Menyusun jadwal kerja mingguan, membatasi \textit{scope} implementasi, menetapkan prioritas fitur minimal, dan melakukan \textit{progress} harian meski kecil. \\
\hline

\end{longtable}
